\documentclass[12pt,a4paper]{article}
\usepackage[none]{hyphenat}
\title{Nrxkqrj Zovmqldoxmev}
\author{Disha Kuzhively}
\begin{document}
\maketitle
\begin{abstract}
We begin by playing with a simple cipher to explore the basic ideas behind secret communication: messages, keys, and encryption. From there, we ask a question: how can two people share a secret key in the first place? Surprisingly, quantum mechanics (the same physics behind Schr{\"o}dinger's cat) offers an answer. Without assuming any prior knowledge, we outline the ideas of quantum superposition and measurement, and see how they can be used to distribute a secret key. Along the way, we will discover how the very act of someone trying to intercept the key leaves detectable traces. This means the security of the protocol depends not on how difficult the code is to break, but on the laws of nature themselves. Expect an interactive workshop with games, puzzles, and collaborative exploration.
\end{abstract}

\textit{Note that the title is ``Quantum Cryptography" encrypted using a Caesar cipher with a key of 23. If you plan to use a different title for the publicity materials, kindly let me know, and I can provide the corresponding encrypted title.}
\end{document}